\documentclass[11pt]{article}
\usepackage[margin=1in]{geometry}
\usepackage{booktabs}
\usepackage{amssymb}
\usepackage{caption}
\usepackage{xcolor}
\usepackage{float}
\usepackage{graphicx}
\usepackage{listings}

\captionsetup{font=small, labelfont=bf, labelsep=period}
\lstdefinestyle{pythonstyle}{
    language=Python,
    basicstyle=\ttfamily\scriptsize,
    keywordstyle=\color{black}\bfseries,
    commentstyle=\color{gray},
    stringstyle=\color{black},
    showstringspaces=false,
    breaklines=true,
    frame=single,
    framerule=0.2pt,
    rulecolor=\color{gray},
    columns=fullflexible
}

\title{Neural Network Prediction of Climate-Related Supply Chain Risk: AER-Style Figures}
\author{}
\date{}

\begin{document}
\maketitle

\noindent
The figures use the original Stata replication data. The prediction label is
\(\mathbb{1}\{\sum_{h=1}^{4} Ter\_q\_qh>0\}\), defined on observations with all
four future termination variables observed. The neural network is a two-hidden-layer
MLP classifier trained only on pre-holdout observations. No simulated outcome data
are used.

\begin{figure}[H]
\centering
\includegraphics[width=0.98\linewidth]{figures/fig1_sample_facts.pdf}
\caption{Prediction sample and climate exposure. Notes: Panel A reports the annual
share of customer-quarter observations with a positive increase in relationship
terminations over the following four quarters. Panel B reports the sample mean of
the customer-city extreme-weather variable and the climate-finance exposure variable.
Light bars flag sparse early-year cells with fewer than 50 observations.}
\end{figure}
\noindent\textbf{Python code for Figure 1.}
\lstinputlisting[style=pythonstyle]{code/figure_1_sample_facts.py}
\clearpage

\begin{figure}[H]
\centering
\includegraphics[width=0.98\linewidth]{figures/fig2_prediction_curves.pdf}
\caption{Out-of-sample predictive performance. Notes: The holdout period is
2022-2023. The logit curve is included only as a transparent benchmark for the
neural network, using the same input variables and train-test split.}
\end{figure}
\noindent\textbf{Python code for Figure 2.}
\lstinputlisting[style=pythonstyle]{code/figure_2_prediction_curves.py}
\clearpage

\begin{figure}[H]
\centering
\includegraphics[width=0.98\linewidth]{figures/fig3_calibration_residuals.pdf}
\caption{Calibration and residual diagnostics. Notes: Predictions are produced
by a rolling one-year-ahead design. Panel A bins neural-network predictions into
deciles. Panel B plots autocorrelations of quarterly mean residuals; the shaded
band is \(\pm 1.96/\sqrt{T}\).}
\end{figure}
\noindent\textbf{Python code for Figure 3.}
\lstinputlisting[style=pythonstyle]{code/figure_3_calibration_residuals.py}
\clearpage

\begin{figure}[H]
\centering
\includegraphics[width=0.98\linewidth]{figures/fig4_model_stability.pdf}
\caption{Model stability diagnostics. Notes: Panel A follows the parameter-distance
logic of the model-stability framework. Distances are normalized Frobenius distances
between MLP parameter vectors and the 2018 reference model. The dashed line is the
95th percentile of same-sample distances generated by changing only the random seed.
Panel B reports the corresponding one-year-ahead AUC.}
\end{figure}
\noindent\textbf{Python code for Figure 4.}
\lstinputlisting[style=pythonstyle]{code/figure_4_model_stability.py}
\clearpage

\begin{figure}[H]
\centering
\includegraphics[width=0.86\linewidth]{figures/fig5_permutation_importance.pdf}
\caption{Permutation importance in the holdout sample. Notes: Each bar is the
average decrease in holdout AUC when the feature is permuted 50 times. Error bars
are 95 percent simulation intervals across permutations.}
\end{figure}
\noindent\textbf{Python code for Figure 5.}
\lstinputlisting[style=pythonstyle]{code/figure_5_permutation_importance.py}
\clearpage

\begin{figure}[H]
\centering
\includegraphics[width=0.98\linewidth]{figures/fig6_risk_surface.pdf}
\caption{Model-implied supply chain risk surfaces. Notes: Curves average predicted
risks over the 2022-2023 holdout covariate distribution. Panel A varies
climate-finance exposure and supplier financing constraints while updating their
interaction exactly as in the original data. Panel B varies customer-city climate
shock intensity; tick marks show the empirical support of the holdout distribution.}
\end{figure}
\noindent\textbf{Python code for Figure 6.}
\lstinputlisting[style=pythonstyle]{code/figure_6_risk_surface.py}
\clearpage

\section*{Shared Python Utilities}
The figure scripts above import the following shared functions for data loading,
target construction, model training, prediction, and parameter-distance diagnostics.
\lstinputlisting[style=pythonstyle]{code/common.py}

\end{document}
